\documentclass[11pt]{article}
\usepackage[utf8]{inputenc}
\usepackage[T1]{fontenc}
\usepackage{natbib}
\usepackage{hyperref}
\usepackage{titling}
\usepackage[left=1in,right=1in,bottom=1in]{geometry}
\usepackage{setspace}
\usepackage{amsmath}

\setlength{\droptitle}{-7em}
\setstretch{1.2}

\title{Literature Summary}
\author{Damanveer Singh Dhaliwal}
\date{\today}

\begin{document}
\begin{titlepage}
    \centering
    \vspace{0.5cm}

    \Large{\textbf{Regime-Switching Gaussian Process Cross-Sectional Asset Pricing with Covariate Driven Transitions}}\\  
    \Large{STA2104 - Proposal}\\  
    \vspace{0.5cm}

    \Large{\textbf{Author Name:}}\\
    \Large{Damanveer Singh Dhaliwal}\\
    \Large{Student ID: 1012787147} \\
    \Large{Email: d.dhaliwal@mail.utoronto.ca}

    \vspace{2cm}
    \textbf{Abstract} \\
    \vspace{0.5cm}
    \begin{minipage}{\textwidth}
    \normalsize
    The relationship between firm characteristics and expected stock returns is both non-linear and time-varying. Yet, most pricing models assume either a fixed functional form or a stationary relationship. To overcome these two drawbacks, we propose a new asset pricing model that combines three components: Gaussian process to non-parametrically learn how the characteristics of the underlying assets affect stock returns, hidden Markov model that allows the pricing function to switch between different latent market states, non-homogeneous transition probabilities that are driven by observable macroeconomic variables such as volatility, credit spread, etc. The model learns regime-dependent importance of underlying asset characteristics, offering an interpretable answer to the question of which stock characteristics matter in which market environment. We validate the approach on synthetic data first with known ground truth, before applying it to the cross-section of actual market data of S\&P 500 returns.
    \end{minipage}

    \vfill
    
    \large{Department of Economics}\\
    \large{University of Toronto}
    
    \vspace{0.8cm}
    
    \large{\today}
\end{titlepage}
\vspace{-2em}
\section{Introduction}
The cross-section of asset returns involves modeling the returns of a large number of assets at a given point in time. The central question in empirical finance is to understand why some stocks earn higher returns than others. This is modeled by the relationship between the expected stock returns and the relationship between firm characteristics such as size, value, momentum, profitability. However, there are two specific challenges that arise when modeling this. 

First, the functional form of this relationship is unknown. \cite{fama_risk_1973} proposed a linear regression model to explain the cross-section of expected returns. More recently, \cite{gu_empirical_2020} showed that non-linearity and interactions among characteristics matter substantially. Second, the relationship is not stationary. Some stocks outperform in some environments and underperform in others. Momentum crashes during market transitions, and factor premium varies with macroeconomic conditions. Most existing approaches, including recent ML methods, estimate a single time-invariant pricing function. 

To address these two problems, I propose a model by combining three components:
\begin{enumerate}
    \item \textbf{Gaussian Process (GP):} Within each regime, a Gaussian process learns the characteristic return mapping non-parametrically, providing both flexibility and calibrated uncertainty. 
    \item \textbf{Hidden Markov Model (HMM): }The HMM allows the pricing function itself to switch between $K$ discrete latent regimes, so that different market environments can have fundamentally different relationships between the characteristics and the returns of the firm. 
    \item \textbf{Non-Homogeneous Transition Probabilities:} The transition probabilities of the HMM governing the state of the world or the regime that we believe the current time period is in is not constant but depends on observable covariates, such as macroeconomic indicators, market volatility, or other relevant factors. This allows our model to learn when the regime shifts are more likely based on current economic conditions rather than relying purely on latent dynamics. 
\end{enumerate}

The three components combine into a single generative model. At each time period $t$, a latent regime state is drawn from the non-homogeneous HMM. Conditional on that regime, a Gaussian process with regime-specific hyperparameters takes the input of the cross-section of the stock characteristics and outputs the expected return. The observed returns are assumed to be noisy realisations of the GP.

Concretely at time $t$, the return of stock $n$ conditional on regime $s_t = k$ is:
\begin{equation*}
    r_{n,t} = f_k(\textbf{x}_{n,t}) + \epsilon_{n,t}, \quad f_k \sim \mathcal{GP}(0, \kappa_k), \quad \epsilon_{n,t} \sim \mathcal{N}(0, \sigma^2_k)
\end{equation*}
where $\textbf{x}_{n,t}$ is the vector of firm characteristics and $\kappa_k$ is the Automatic Relevance Determination (ARD) kernel function for regime $k$. The regime transitions follow a logistic model of observable macroeconomic covariates.

Inference uses a two-stage Expectation-Maximization (EM) algorithm inspired by \cite{li_hidden_2023} and incorporates covariate-driven transitions. The E-step runs forward-backward on the HMM using the GP marginal likelihood $p(r_t | s_t = k) = \mathcal{N}(r_t | 0, K_k + \sigma^2_k I)$ as the emission probability. Here, the function $f$ is marginalised out at each time step. The regime defines a distribution over the pricing functions through its hyperparameters, and this preserves the conditional independence required by the forward-backward algorithm of HMM. 

The M-step updates parameters by weighting each time period by how likely it is to belong to each regime. Time periods that are more likely in regime $k$ have more influence on regime $k$'s GP hyper-parameters. Transition coefficients are updated the same way, weighted by how likely each pair of consecutive regimes is. 


A summary of the literature has been presented in section 4 but initial search confirmed that no existing work has combined all three components in cross-sectional asset pricing.

\section{Work Plan}
The following is the  proposed work plan to implement the model and evaluate its performance. To satisfy the fail-fast principle, I will validate the feasibility on synthetic data before committing to real-data infrastructure.
\begin{enumerate}
    \item Extend literature search and finalize the related work section. This will involve a more comprehensive review of existing models in asset pricing, Gaussian processes, and regime-switching models to ensure that the proposed approach is novel and to identify any potential pitfalls or areas for improvement (3 hours).
    \item Implement the GP regression with Automatic Relevance Determination (ARD) kernel and marginal likelihood optimization (3 hours).
    \item Toy Check 1: Generate data from a known GP with 2 relevant and 3 irrelevant dimensions and verify parameter recovery and posterior coverage (1 hour).
    \item Implement the HMM component and the EM algorithm for inference (2 hours).
    \item Toy Check 2: Simulate a 2-state HMM with known parameters and verify that the model can recover the transition matrix and emission parameters (1 hour).
    \item Implement the non-homogeneous transition model (4 hours).
    \item Toy Check 3: Simulate HMM with covariate-driven transitions and verify that the model can recover the transition coefficients (1 hour).
    \item Integrate all three components into a single pipeline (3 hour)
    \item Final Toy Check: Generate 100 months of 100 stock cross sections from a 2-regime model with distinct GPs and macro-driven transitions. Verify regime identification, lengthscale recovery and predictive improvement over single-GP baseline (1 hour).
    \item Download and preprocess S\&P 500 returns, firm characteristics and macroeconomic data (0.5 hours).
    \item Run full model on real data with $K=2$ and compute out-of-sample predictive performance and interpret the discovered regimes (1 hour).
    \item Implement and run baseline models: \cite{fama_risk_1973}, single-regime GP, random forest (2 hours).
    \item Generate all figures, tables, and regime analysis plots (2 hours).
    \item Write up the results and discussion sections (5 hours).
\end{enumerate}
\section{Proposed Results}
The proposed results of this project are expected to demonstrate the effectiveness of the regime-switching Gaussian process model in capturing the non-linear and time-varying relationships between firm characteristics and stock returns. The key results will include:
\begin{enumerate}
    \item Figure 1: We use toy data to recover true vs learned ARD length scales across both regimes on synthetic data. This figure informs the reader whether the inference algorithm can identify which characteristics matter in each regime when the ground truth is known. 
    \item Figure 2: We use toy data now to infer regime probabilities overlaid on the true regime sequence for the synthetic dataset. This informs the reader on how accurately and how quickly the HMM detects regime switches. 
    \item Figure 3: We estimate posterior regime probabilities on real data. This figure will inform the reader whether the inferred regimes correspond to economically meaningful periods that would be overlaid on the graph, such as the dot-com bust, the global financial crisis, the COVID-19 pandemic. 
    \item Figure 4: We estimate a regime-dependent characteristics that are important. We create a bar chart of the length scales for each characteristic separated by regime. This informs the reader which stock characteristics the model considers most important in each market environment and will become the key interpretability output of this project. 
    \item Figure 5: We test our assumptions and do a side-by-side comparison of out-of-sample performance under different kernels. This informs us and the readers the sensitivity of our algorithm to the specific assumptions that we make. 
    \item Table 1: We compare the out-of-sample performance of our model against the baselines. We use a monthly cross-sectional R-squared and try to inform the reader whether the regime switching and non-parametric flexibility together improve prediction over existing alternatives. 
    \item Table 2: To analyze economic interpretation, we compare annualised Sharpe ratios of long-short portfolios based on each model's predicted returns. This informs us whether better statistical predictions do translate into the real world and meaningful improvement in portfolio management.
\end{enumerate}

\section{Related Work}
\cite{filipovic_empirical_2025} is the closest paper to this proposal. They also apply a GP regression to the cross-section of stock returns, offering non-parametric flexibility with principled uncertainty quantification. Unlike their model, which estimates a single time-invariant pricing function, we allow the GP to switch across regimes, capturing the non-stationarity of the underlying relationship between stock returns and asset characteristics. \cite{guidolin_size_2008} show that a regime-switching linear model outperforms a single-regime linear model.

\cite{gu_empirical_2020} applies machine learning methods such as neural nets, random forests, and gradient boosting trees to show that they significantly outperform linear models for cross-sectional return prediction, establishing that non-linearity and interactions with the characteristics of the underlying asset offer better predictions. Unlike their work, we provide a Bayesian framework with regime structure and calibrated uncertainty rather than a single black box prediction function. 

\cite{jung_scalable_2022} develop scalable inference for a Bayesian hidden Markov model with GP emission models where each hidden state has its own GP with state-specific kernel hyperparameters. We inherit the core idea of regime-specific GPs governed by an HMM from them. Our work differs from theirs in two places: they target time series clustering with one-dimensional scalar inputs, while we apply it to cross-sectional asset pricing with high-dimensional characteristic vectors. They also use standard homogeneous transition probabilities, whereas to incorporate economic theory, we introduce a covariate-driven transition probability that varies.

\cite{li_hidden_2023} also apply a regime-switching GP model to time series data but they focus on electricity load forecasting. Our inference approach is inspired by theirs, where we fit HMM parameters first, then estimate per-regime GP parameters. However, unlike their work which operates on a univariate time series with seasonal structure, our setting involves a cross-section of hundreds of stocks with characteristic-driven pricing, which requires us to interpret which features are more relevant and to also model non-homogeneous transitions to capture macroeconomic regime shifts. 



\pagebreak
\bibliographystyle{chicago}
\bibliography{references}

\end{document}